% Oblikovanje po Wordovi predlogi FS, september 2026.
\usepackage{iftex}
\ifPDFTeX
  \PackageError{FS2026}{Uporabite XeLaTeX ali LuaLaTeX}{Pisava Times New Roman zahteva XeLaTeX ali LuaLaTeX.}
\fi
\usepackage{fontspec}
\setmainfont{Times New Roman}
\setmonofont{Latin Modern Mono}[Scale=MatchLowercase]
\usepackage[english,slovene]{babel}
\usepackage{amsmath,amssymb,bm}
\usepackage{graphicx}
\usepackage[section]{placeins}
\usepackage{float}
\usepackage{enumitem}
\usepackage{longtable,booktabs,array}
\usepackage{fancyhdr}
\usepackage{tocloft}
\usepackage{titlesec}
\usepackage{pdfpages}
\usepackage{pdflscape}
\usepackage{caption,subcaption}
\usepackage{icomma}
\usepackage{etoolbox}
\usepackage{notoccite}
\usepackage[numbers,square,sort&compress]{natbib}
\usepackage[inner=30mm,outer=25mm,top=30mm,bottom=25mm,headheight=15pt]{geometry}
\usepackage[unicode,hidelinks]{hyperref}

% Umerjeno po dejanskem razmiku osnovnic v Wordovem PDF-ju (pribl. 15,9 bp).
% Times New Roman ima v XeLaTeXu drugačne metrike, zato faktor ni neposredno 1,15.
\linespread{1.10}
\setlength{\parindent}{0pt}
% Sosednja odmika Wordovih odstavkov po 8 pt se ne seštevata.
\setlength{\parskip}{8pt}
\raggedbottom
\setlength{\emergencystretch}{1em}
\setlength{\mathindent}{5mm}
\setcounter{secnumdepth}{2}
\setcounter{tocdepth}{2}
\setlist[itemize,1]{label=\textbullet,leftmargin=7.5mm,itemsep=2pt,parsep=0pt,topsep=2pt}
\setlist[itemize,2]{label=\(\circ\),leftmargin=7.5mm,itemsep=8pt,parsep=0pt,topsep=2pt}
\setlist[enumerate]{leftmargin=*,itemsep=2pt,parsep=0pt}
\DeclareCaptionFont{fsfont}{\fontsize{11}{12.65}\selectfont}
\captionsetup{font=fsfont,labelsep=colon}
\captionsetup[table]{position=top,justification=raggedright,singlelinecheck=false,skip=4pt}
\captionsetup[figure]{position=bottom,justification=centering,skip=4pt}
\captionsetup[subfigure]{labelfont=normalfont}
% \@floatboxreset znotraj plavajočega okolja kliče \normalsize, zato velikost
% pisave v preglednicah nastavimo za njim (\AtBeginEnvironment{table} ne deluje)
\makeatletter
\appto\@floatboxreset{\ifdefstring{\@captype}{table}{\fontsize{10}{11.5}\selectfont}{}}
\makeatother
\addto\captionsslovene{\renewcommand{\tablename}{Preglednica}\renewcommand{\bibname}{Literatura}}
\setlength{\textfloatsep}{12pt}
\setlength{\intextsep}{12pt}
\newcommand{\minitab}[2][l]{\begin{tabular}{#1}#2\end{tabular}}
\renewcommand{\thetable}{\thechapter.\arabic{table}}
\renewcommand{\thefigure}{\thechapter.\arabic{figure}}
\renewcommand{\theequation}{\thechapter.\arabic{equation}}

\titleformat{\chapter}[hang]{\fontsize{24}{27.6}\selectfont\bfseries}{\thechapter}{1em}{}
\titleformat{name=\chapter,numberless}[hang]{\fontsize{24}{27.6}\selectfont\bfseries}{}{0pt}{}
% Odmiki so umerjeni po prvi oštevilčeni strani Wordove predloge.
\titlespacing*{\chapter}{0pt}{111pt}{36pt}
\titleformat{\section}{\fontsize{18}{20.7}\selectfont\bfseries}{\thesection}{1em}{}
\titlespacing*{\section}{0pt}{15.25pt}{6pt}
\titleformat{\subsection}{\fontsize{14}{16.1}\selectfont\bfseries}{\thesubsection}{1em}{}
\titlespacing*{\subsection}{0pt}{18pt}{6pt}
\titleformat{\subsubsection}{\normalsize\bfseries}{}{0pt}{}
\titlespacing*{\subsubsection}{0pt}{12pt}{3pt}
\renewcommand{\cftchapdotsep}{\cftdotsep}
% Naslovi prve ravni (\chapter) so v kazalu veliki 16 pt.
\renewcommand{\cftchapfont}{\fontsize{16}{19.27}\selectfont\bfseries}
\renewcommand{\cftchappagefont}{\fontsize{16}{19.27}\selectfont\bfseries}
% Naslovi druge ravni (\section) so v kazalu izpisani krepko.
\renewcommand{\cftsecfont}{\bfseries}
\renewcommand{\cftsecpagefont}{\bfseries}
% Navpični razmiki v kazalu so umerjeni po Wordovi predlogi (razmiki osnovnic:
% poglavje-poglavje 29,2 pt, poglavje-razdelek 27,5 pt, razdelek-razdelek 26 pt,
% razdelek-podrazdelek 21,8 pt). Odstavek vpisa poglavja se konča v njegovi
% pisavi, da velja razmik vrstic 16-pt pisave; kern za vpisom ne vpliva na \addvspace.
\setlength{\cftbeforechapskip}{6.3pt}
\setlength{\cftbeforesecskip}{10.1pt}
\setlength{\cftbeforesubsecskip}{5.9pt}
\renewcommand{\cftchapafterpnum}{\fontsize{16}{19.27}\selectfont\par\kern1.7pt}

\fancypagestyle{fancy_nohead}{%
  \fancyhf{}\fancyfoot[LE,RO]{\thepage}
  \renewcommand{\headrulewidth}{0pt}\renewcommand{\footrulewidth}{0pt}}
\fancypagestyle{plain}{%
  \fancyhf{}\fancyfoot[LE,RO]{\thepage}
  \renewcommand{\headrulewidth}{0pt}}
\renewcommand{\chaptermark}[1]{\markboth{#1}{}}
\fancypagestyle{vsebina}{%
  \fancyhf{}
  \fancyhead[LE,RO]{\nouppercase{\leftmark}}
  \fancyfoot[LE,RO]{\thepage}
  \renewcommand{\headrulewidth}{0.1mm}
  \renewcommand{\footrulewidth}{0pt}}

% Povsem prazna stran brez številke in glave (npr. vezni list).
\newcommand{\praznastran}{\null\thispagestyle{empty}\newpage}
\makeatletter
% Soda prazna stran pred novim poglavjem ima, kot v Wordu, številko in glavo
% trenutnega sloga strani (na naslovnih straneh je slog empty).
\renewcommand{\cleardoublepage}{\clearpage\if@twoside\ifodd\c@page\else\null\newpage\fi\fi}
\newcommand{\vsebinskokazalo}{\begingroup\setlength{\parskip}{0pt}\@starttoc{toc}\endgroup}
\makeatother
% Združljivost s starimi poglavji; novih ročnih praznih strani ni treba dodajati.
\newcommand{\addifeven}{\cleardoublepage}

% Formalni naslovi na vrhu strani; samo seznama se vpišeta v kazalo.
\newcommand{\predpoglavje}[1]{%
  \cleardoublepage\thispagestyle{fancy_nohead}\markboth{#1}{}
  % Umerjeno po Wordu: osnovnica naslova 104 pt od vrha strani, zgornji rob
  % črte 7,8 pt pod osnovnico (neodvisno od spodnjih zank črk), prva vrstica
  % besedila 35,6 pt pod črto.
  \vspace*{-23.6pt}{\fontsize{18}{20.7}\selectfont\bfseries #1\par}
  \nobreak\vskip-\prevdepth\vskip7.6pt\hrule height 0.4pt\nobreak\vspace{19.2pt}}
\newcommand{\predpoglavjevkazalu}[1]{%
  \predpoglavje{#1}\phantomsection\addcontentsline{toc}{chapter}{\protect\brezzamika #1}}
% Neoštevilčeni vpisi v kazalu nimajo visečega zamika druge vrstice (kot v Wordu).
% Ukaz ponastavi levi zamik (globalno, ker je vpis v skupini; privzeta vrednost
% je tako in tako 0) in izravna začetni odmik tocloft.
\newcommand{\brezzamika}{\global\leftskip=\cftchapindent\relax\hskip\cftchapnumwidth\relax}
\AtBeginDocument{\pdfstringdefDisableCommands{\let\brezzamika\relax}}
\newcounter{priloga}
\renewcommand{\thepriloga}{\Alph{priloga}}
\newcommand{\priloga}[1]{%
  \refstepcounter{priloga}
  \chapter*{Priloga \Alph{priloga} -- #1}
  \phantomsection
  \addcontentsline{toc}{chapter}{\protect\brezzamika Priloga \Alph{priloga} -- #1}
  \markboth{Priloga \Alph{priloga}}{}
  \renewcommand{\thechapter}{\Alph{priloga}}
  \setcounter{figure}{0}\setcounter{table}{0}\setcounter{equation}{0}}
\setlength{\bibsep}{8pt}
\renewcommand{\bibfont}{\normalfont\normalsize}
\urlstyle{same}
\AtBeginDocument{\hypersetup{pdftitle={\naslovSlo},pdfauthor={\avtor},pdfsubject={\tipDela}}}

\newcommand{\poglavjeref}[1]{%
	\textbf{\ref{#1}~\nameref{#1}}%
}
